\chapter{Evaluation Criteria}

\vspace*{1cm}

The evaluation of a fine-tuned language model should not be limited to the final
training loss. Since this project targets Edge AI applications, the comparison
must also include efficiency-related criteria. A model that gives good prediction
quality but requires high memory, long inference time, or excessive energy may
not be suitable for deployment on constrained devices.

\section{Validation Loss}

Validation loss measures how well the model performs on data that was not used
for parameter updates during training. It is one of the main indicators of model
generalization. A lower validation loss usually means that the model is better at
predicting the expected outputs on unseen examples.

\vspace*{0.5cm}

In this project, validation loss is used to compare the fine-tuning quality of
the different optimizers. It helps identify whether an optimizer improves the
model or causes unstable training.

\section{Perplexity}

Perplexity is a common metric in language modeling. It is derived from the loss
and gives an interpretation of how uncertain the model is when predicting the next
tokens. Lower perplexity indicates better language modeling performance.

\vspace*{0.5cm}

Perplexity can be computed from the validation loss as:

\begin{equation}
    \text{Perplexity} = e^{\text{Validation Loss}}
\end{equation}

This metric is useful because it gives another view of model quality. However, it
must be interpreted together with efficiency metrics when the target environment
is Edge AI.

\section{Training Time}

Training time represents the total time required to fine-tune the model for a
given optimizer. This metric is important because some algorithms may produce good
results but require much longer execution time.

\vspace*{0.5cm}

For example, L-BFGS may require several evaluations of the loss during one update
step, which can increase training time compared with optimizers such as Adam. In
an experimental comparison, training time helps evaluate the practical cost of
each optimizer.

\section{Training Speed}

Training speed is measured as the number of processed examples per second. It
provides a normalized view of how fast the training process runs.

\vspace*{0.5cm}

This metric is useful when comparing optimizers because it shows how efficiently
each method uses the available hardware. A higher training speed is generally
better, especially when experiments must be repeated many times.

\section{Memory Consumption}

Memory consumption measures the amount of memory used during training. For GPU
training, this usually includes the allocated and reserved GPU memory. This
criterion is very important for Edge AI because target devices often have limited
available memory.

\vspace*{0.5cm}

An optimizer may require additional memory to store internal states. For example,
Adam stores moving averages of gradients and squared gradients. L-BFGS also keeps
a history of previous updates. Therefore, memory usage can differ significantly
between optimizers.

\section{Model Size}

Model size represents the storage space required by the final fine-tuned model.
Smaller models are easier to deploy on edge devices because they require less disk
space and can be loaded more easily into memory.

\vspace*{0.5cm}

In this project, model size is measured after saving the fine-tuned model. It is
used as a baseline for future compression experiments such as quantization or
pruning.

\section{Inference Time and Latency}

Inference time measures the time required by the model to generate outputs for a
set of examples. Latency represents the average response time per sample. These
two metrics are essential for real-time or interactive Edge AI applications.

\vspace*{0.5cm}

Low latency is especially important when the model must respond quickly to user
requests. Even if two optimizers produce similar validation loss, the one that
leads to faster inference can be more suitable for practical deployment.

\section{Energy Consumption}

Energy consumption is another important criterion for Edge AI. Devices such as
mobile phones, sensors, and embedded boards often operate with limited battery
capacity. A model that consumes less energy is more appropriate for long-term or
portable use.

\vspace*{0.5cm}

In this project, energy consumption can be estimated when NVIDIA power telemetry
is available. If telemetry is not available, the metric can be left unavailable
and discussed as a limitation.

\section{Compression and Performance Tradeoff}

The final criterion is the compression and performance tradeoff. The fine-tuned
models produced by this project are considered baselines for future compression
experiments. Compression techniques aim to reduce model size and computational
cost while preserving acceptable performance.

\vspace*{0.5cm}

The optimizer comparison therefore helps determine which fine-tuned model is the
best starting point for later compression. A good candidate should combine strong
validation performance with reasonable memory usage, latency, and model size.

\section{Summary of Metrics}

The main evaluation criteria are summarized in Table~\ref{tab:evaluation_metrics}.

\begin{table}[H]
    \centering
    \begin{tabularx}{\textwidth}{lX}
        \toprule
        \textbf{Metric} & \textbf{Purpose} \\
        \midrule
        Validation loss & Measures model generalization on unseen data. \\
        Perplexity & Evaluates language modeling uncertainty. \\
        Training time & Measures total optimization cost. \\
        Training speed & Measures examples processed per second. \\
        Memory consumption & Evaluates hardware resource usage. \\
        Model size & Measures storage cost of the fine-tuned model. \\
        Inference time & Measures total generation time. \\
        Latency & Measures average response time per sample. \\
        Energy consumption & Estimates power efficiency when telemetry is available. \\
        Compression tradeoff & Prepares comparison with future compressed models. \\
        \bottomrule
    \end{tabularx}
    \caption{Evaluation criteria used to compare the optimizers.}
    \label{tab:evaluation_metrics}
\end{table}
